%! AP Statistics — Unit 9, Lesson 9.6 — Student Packet Cover Sheet
\documentclass[10pt]{article}
\usepackage{apstats-article}
\usepackage{apstats-boxes}
\usepackage{ltablex}
\keepXColumns
\usepackage{tikz}

\begin{document}

\begin{tikzpicture}[remember picture, overlay]
  \fill[navy] (current page.north west)
    rectangle ([yshift=-0.9in]current page.north east);
\end{tikzpicture}
\vspace{-0.8in}

\begin{center}
  {\color{white}\textbf{\LARGE Statistics}}\\[4pt]
  {\color{skymid}\large\textbf{Unit 9}}\\[6pt]
  {\color{white}\textbf{\large Lesson 9.6 \quad Skills Focus: Selecting an Appropriate Inference Procedure}}
\end{center}

\vspace{0.22in}
\namedateperiod
\vspace{0.18in}

\begin{learningtargetbox}
\small
\begin{enumerate}[leftmargin=1.5em, itemsep=5pt]
  \item I can identify whether a context calls for a confidence interval or a significance test.
  \item I can select the correct inference procedure for a given context from among all
        procedures in Units 6--9.
  \item I can state $H_0$ and $H_a$ (or describe the parameter being estimated) for the
        selected procedure.
\end{enumerate}
\end{learningtargetbox}

\vspace{0.14in}

\begin{tocbox}
% A CONTENTS LIST ONLY: no Score column and no Total row. Whether a component counts is said in
% its description.
\small
\renewcommand{\arraystretch}{1.5}
\begin{tabularx}{\linewidth}{@{} c l X @{}}
\rowcolor{sky}
\textbf{\#} & \textbf{Component} & \textbf{Description} \\
\midrule
1 & Warm-Up        & Classify three scenarios: variable type, group structure, procedure \\
2 & Guided Notes   & Decision framework; reference table; quick practice \\
3 & Group Activity & 8-scenario identification: classify, name procedure, state parameter \\
4 & Exit Ticket    & Two scenarios: name procedure and justify \\
5 & Homework       & Five AP-style scenarios: procedure, parameter/hypotheses, conditions \\
\end{tabularx}
\end{tocbox}

\end{document}
